\section{m09 Boosting}

% source: wiki/mocs/m09-boosting.md:17 — "weak learners + a loss + an additive way to combine them"
% source: archive/modules/9TreeBoosting/9TreeBoosting.Rmd:L221 — definition of boosting
Boosting: iteratively add basis functions (typically shallow trees) in a greedy fashion so each addition reduces the chosen loss. Applied to models with \emph{high bias / low variance} (esp.\ trees). Final model is an additive expansion
\[ f(x)=\sum_{m=1}^M \beta_m\, b(x;\gamma_m), \qquad f_M(x)=\sum_{m=1}^M T(x;\Theta_m). \]
% source: archive/modules/9TreeBoosting/9TreeBoosting.Rmd:L550-556, L614-616

\subsection{Three main ingredients}
% source: archive/modules/9TreeBoosting/9TreeBoosting.Rmd:L908-919
\begin{tabular}{@{}p{0.22\linewidth} p{0.72\linewidth}@{}}
\toprule
Ingredient & Role \\
\midrule
Weak learner & High-bias / low-variance base (typ.\ shallow tree / stump). \\
Loss $L$ & Differentiable cost we want to minimise. \\
Additive model & Sequentially add weak learners, each lowering $L$. \\
\bottomrule
\end{tabular}

\subsection{AdaBoost.M1 (Algorithm 10.1, ESL)}
% source: archive/modules/9TreeBoosting/9TreeBoosting.Rmd:L282-301
Binary classification, $Y\in\{-1,+1\}$, weak classifier $G(x)$.

\begin{algorithm}[H]
\caption{AdaBoost.M1 (ESL Alg.\ 10.1)}
\begin{algorithmic}[1]
\State Initialise weights $w_i = 1/N$, $i=1,\dots,N$.
\For{$m = 1,2,\dots,M$}
  \State Fit $G_m(x)$ to training data using weights $w_i$.
  \State $\text{err}_m \gets \dfrac{\sum_{i=1}^N w_i\, \mathbb{I}(y_i \neq G_m(x_i))}{\sum_{i=1}^N w_i}$
  \State $\alpha_m \gets \log\!\big((1-\text{err}_m)/\text{err}_m\big)$
  \State $w_i \gets w_i \cdot \exp\!\big(\alpha_m\, \mathbb{I}(y_i \neq G_m(x_i))\big)$, $i=1,\dots,N$
\EndFor
\State \Return $G(x) = \operatorname{sign}\!\Big[\sum_{m=1}^M \alpha_m G_m(x)\Big]$
\end{algorithmic}
\end{algorithm}

% source: archive/modules/9TreeBoosting/9TreeBoosting.Rmd:L305-317
\textbf{Notes.} $\alpha_m>0$ iff $\text{err}_m<0.5$; misclassified points get multiplied by $\exp(\alpha_m)$, correctly classified ones keep their weight; output is the sign of a continuous score. AdaBoost equals forward-stagewise additive modelling with the \emph{exponential loss} $L(y,f)=\exp(-yf)$ (discovered 5 yrs after invention).
% source: archive/modules/9TreeBoosting/9TreeBoosting.Rmd:L633-642
(See ESL Fig.\ 10.1 for the schematic.)

\subsection{Boosting regression trees (ISLP Alg.\ 8.2)}
% source: archive/modules/9TreeBoosting/9TreeBoosting.Rmd:L428-442
\begin{algorithm}[H]
\caption{Boosting for regression trees (ISLP Alg.\ 8.2)}
\begin{algorithmic}[1]
\State Set $\hat f(x)=0$ and $r_i=y_i$ for all $i$ in training set.
\For{$b=1,2,\dots,B$}
  \State Fit tree $\hat f^{b}$ with $d$ splits ($d{+}1$ terminal nodes) to data $(X,r)$.
  \State Update $\hat f(x) \gets \hat f(x) + \nu\, \hat f^{b}(x)$
  \State Update residuals $r_i \gets r_i - \nu\, \hat f^{b}(x_i)$
\EndFor
\State \Return $\hat f(x) = \sum_{b=1}^B \nu\, \hat f^{b}(x)$
\end{algorithmic}
\end{algorithm}
Three tunings: $B$ (number of trees), $\nu$ (shrinkage), $d$ (interaction depth). $d=1$ gives an additive model.
% source: archive/modules/9TreeBoosting/9TreeBoosting.Rmd:L442-458

\subsection{Forward-stagewise additive modelling (boosting trees)}
% source: archive/modules/9TreeBoosting/9TreeBoosting.Rmd:L564-577
Single tree as $T(x;\Theta)=\sum_{j=1}^J \gamma_j\, \mathbb{I}(x\in R_j)$, parameters $\Theta=\{R_j,\gamma_j\}_1^J$. Boosting solves at step $m$:
\[ \hat\Theta_m = \arg\min_{\Theta_m} \sum_{i=1}^N L\!\big(y_i,\, f_{m-1}(x_i) + T(x_i;\Theta_m)\big). \]
% source: archive/modules/9TreeBoosting/9TreeBoosting.Rmd:L621-627
\textbf{Two sub-problems:} (i) given $R_j$, find $\gamma_j$ (easy: mean / majority vote); (ii) find $R_j$ (hard, greedy / surrogate loss).
% source: archive/modules/9TreeBoosting/9TreeBoosting.Rmd:L584-596

\subsection{Steepest descent in function space}
% source: archive/modules/9TreeBoosting/9TreeBoosting.Rmd:L666-686
Loss $L(f)=\sum_i L(y_i,f(x_i))$ has gradient vector $\mathbf{g}_m^\top=(g_{1m},\dots,g_{Nm})$ with components
\[ g_{im}=\Big[\tfrac{\partial L(y_i,f(x_i))}{\partial f(x_i)}\Big]_{f=f_{m-1}}. \]
Update via line-search step $\rho_m = \arg\min_\rho L(\mathbf{f}_{m-1}-\rho\mathbf{g}_m)$, then $\mathbf{f}_m = \mathbf{f}_{m-1} - \rho_m \mathbf{g}_m$.
% source: archive/modules/9TreeBoosting/9TreeBoosting.Rmd:L696-706
(See ESL Fig.\ for steepest descent; \texttt{graphics/Steepest\_descet.png}.)

\subsection{Gradient tree boosting (ESL Alg.\ 10.3)}
% source: archive/modules/9TreeBoosting/9TreeBoosting.Rmd:L760-789
Central idea: fit a tree as close as possible to the \emph{negative gradient}:
\[ \tilde\Theta_m = \arg\min_{\Theta} \sum_{i=1}^N (-g_{im} - T(x_i;\Theta))^2. \]
% source: archive/modules/9TreeBoosting/9TreeBoosting.Rmd:L724-730
That is: fit tree to negative gradient by least squares; then optimise $\gamma_{jm}$ on the resulting regions.

\begin{algorithm}[H]
\caption{Gradient tree boosting (ESL Alg.\ 10.3)}
\begin{algorithmic}[1]
\State Initialise $f_0(x) = \arg\min_\gamma \sum_{i=1}^N L(y_i,\gamma)$.
\For{$m=1$ to $M$}
  \State For $i=1,\dots,N$ compute pseudo-residuals
    \[ r_{im} = -\Big[\tfrac{\partial L(y_i,f(x_i))}{\partial f(x_i)}\Big]_{f=f_{m-1}} \]
  \State Fit regression tree to targets $r_{im}$ giving regions $R_{jm}$, $j=1,\dots,J_m$.
  \State For $j=1,\dots,J_m$ compute
    \[ \gamma_{jm} = \arg\min_\gamma \sum_{x_i\in R_{jm}} L\!\big(y_i, f_{m-1}(x_i) + \gamma\big) \]
  \State Update $f_m(x) = f_{m-1}(x) + \sum_{j=1}^{J_m} \gamma_{jm}\, \mathbb{I}(x\in R_{jm})$
\EndFor
\State \Return $\hat f(x) = f_M(x)$
\end{algorithmic}
\end{algorithm}
(See ESL Fig.\ 10.3 \texttt{graphics/Algorithm10.3.png}.) The $r_{im}$ are called \emph{generalised / pseudo-residuals}.
% source: archive/modules/9TreeBoosting/9TreeBoosting.Rmd:L793-815

\subsection{Loss / objective table}
% source: archive/modules/9TreeBoosting/9TreeBoosting.Rmd:L819-901; wiki/concepts/boosting-loss-functions.md
\scriptsize
\begin{tabularx}{\linewidth}{@{}l X l@{}}
\toprule
Loss & $L(y,f)$ & $-g_i = -\partial L/\partial f$ \\
\midrule
Quadratic & $\tfrac12 (y-f)^2$ & $y-f$ (residual) \\
Absolute & $|y-f|$ & $\operatorname{sign}(y-f)$ \\
Huber & $\begin{cases}(y-f)^2 & |y-f|\le \delta\\ 2\delta|y-f|-\delta^2 & \text{else}\end{cases}$ & piecewise \\
Exponential (AdaBoost) & $\exp(-y f)$, $y\in\{-1,+1\}$ & $y\,\exp(-y f)$ \\
Binomial deviance & $-\mathbb{I}(y{=}1)\log p - \mathbb{I}(y{=}0)\log(1-p)$ & $y - p(x)$ \\
Multinomial deviance & $-\sum_{k=1}^K \mathbb{I}(y{=}k)\log p_k(x)$ & $\mathbb{I}(y_i{=}k) - p_k(x_i)$ \\
\bottomrule
\end{tabularx}
\normalsize

For binomial deviance: $p(x) = e^{f(x)}/(1+e^{f(x)})$ (logistic / sigmoid), $y\in\{0,1\}$.
% source: archive/modules/9TreeBoosting/9TreeBoosting.Rmd:L853-858

For multinomial deviance: $p_k(x) = e^{f_k(x)} / \sum_{l=1}^K e^{f_l(x)}$ (softmax). \textbf{Builds $K$ trees per boosting round}, one per class, each fit to $\mathbf{g}_{km}$; aggregate class probabilities in step 3.
% source: archive/modules/9TreeBoosting/9TreeBoosting.Rmd:L869-899

(See ESL Fig.\ for robust losses, \texttt{graphics/robust\_loss.png}.)

\subsection{Shrinkage / learning-rate update}
% source: archive/modules/9TreeBoosting/9TreeBoosting.Rmd:L1048-1066
Scale each new tree's contribution by $\nu \in (0,1]$:
\[ f_m(x) = f_{m-1}(x) + \nu \sum_{j=1}^{J_m} \gamma_{jm}\, \mathbb{I}(x\in R_{jm}). \]
$\nu$ and $M$ are linked: smaller $\nu \Rightarrow$ larger $M$. Empirically $\nu < 0.1$ is more robust. Strategy: fix $\nu$, choose $M$ by early stopping (validation curve / CV).
% source: archive/modules/9TreeBoosting/9TreeBoosting.Rmd:L1022-1066

\subsection{Boosting-specific hyperparameters}
% source: archive/modules/9TreeBoosting/9TreeBoosting.Rmd:L924-948, L1161-1213; wiki/concepts/xgboost.md
\scriptsize
\begin{tabularx}{\linewidth}{@{}l l X@{}}
\toprule
Symbol & Name & Role \\
\midrule
$M$ (or $B$)         & \# trees / iterations  & too small $\to$ underfit, too large $\to$ overfit; pick via early stopping or CV. \\
$\nu$                & shrinkage / LR         & scales each tree; typ.\ $0.1$ or $0.01$. \\
$d$                  & tree depth (\# splits) & $d{=}1$ stumps $\to$ additive model; common $4\le J\le 8$. \\
$J_m$                & leaves per tree        & $=d+1$; reflects interaction order. \\
$n_{\min}$           & min.\ obs.\ per leaf   & weak knob given $d$ and $M$ are set. \\
$\eta_{\text{row}}$  & row subsample fraction & stochastic GBM; typ.\ $0.5$; \emph{variance} reduction. \\
$\eta_{\text{col}}$  & column subsample fraction & XGBoost variant. \\
$\gamma$ (XGB)       & per-leaf penalty       & post-fit pruning; larger $\gamma\to$ smaller tree. \\
$\lambda$ (XGB)      & L2 on leaf weights     & ridge-style. \\
$\alpha$ (XGB)       & L1 on leaf weights     & lasso-style. \\
$p_{\text{drop}}$ (XGB) & random tree drop    & avoid early-tree dominance (Hinton). \\
\bottomrule
\end{tabularx}
\normalsize

\subsection{Stochastic gradient boosting}
% source: archive/modules/9TreeBoosting/9TreeBoosting.Rmd:L1161-1213
Subsample (without replacement) a fraction $\eta$ of training rows before fitting each tree (Friedman 2002). Typical $\eta=N/2$. Reduces variance \emph{and} computation. Variants: subsample rows, subsample columns per tree, subsample columns per split.

\subsection{XGBoost (eXtreme Gradient Boosting)}
% source: archive/modules/9TreeBoosting/9TreeBoosting.Rmd:L1253-1315; wiki/concepts/xgboost.md
Same forward-stagewise skeleton plus engineering:
\begin{itemize}
\item \textbf{Second-order gradients} (Newton): uses gradient + Hessian (2nd-order Taylor) instead of just gradient.
\item \textbf{Parallelisation}: split-search across CPU cores.
\item \textbf{Approximate split search}: bin features, search bins (not every threshold).
\item \textbf{L1/L2 leaf-weight regularisation} + per-leaf pruning $\gamma$:
\[ \Omega(f) = \gamma |T| + \tfrac12 \lambda \|w\|^2 \quad (\text{L2}), \]
\[ \Omega(f) = \gamma |T| + \alpha \|w\| \quad (\text{L1}). \]
$|T|$ = number of leaves, $w$ = leaf-weight vector.
\item \textbf{Row \& column subsampling} ($\eta_{\text{row}}, \eta_{\text{col}}$).
\item \textbf{Dropout} $p_{\text{drop}}$ (Hinton's trick).
\item \textbf{Shrinkage $\nu$} as in vanilla GBM.
\end{itemize}
Internals (Taylor expansion derivation, histogram split search) are \emph{out of scope}.
% source: wiki/mocs/m09-boosting.md:46-49 — explicit out-of-scope list

\subsection{Partial dependence plots (interpretability)}
% source: archive/modules/9TreeBoosting/9TreeBoosting.Rmd:L1444-1481
For input $X_j$, partial dependence of $f(X)$ on $X_j$:
\[ \E_{X_{-j}} f(X_j, X_{-j}). \]
Estimated for boosted trees by averaging over the data:
\[ \overline{f}(X_j) = \tfrac{1}{N} \sum_{i=1}^N f\!\big(X_j, (X_{-j})_i\big). \]
Effect of $X_j$ \emph{after} accounting for the others; not marginal.

\subsection{Boosted-tree squared-error special case}
% source: archive/modules/9TreeBoosting/9TreeBoosting.Rmd:L633-638, L800-815
For regression tree + squared-error loss, the boosting subproblem is identical to fitting a single regression tree to the previous-step residuals $r_i = y_i - f(x_i)$ — i.e.\ gradient boosting = steepest descent on quadratic loss.

\subsection{Out-of-scope flags (m09)}
% source: wiki/mocs/m09-boosting.md:45-49
Detailed boosting pseudocode line-by-line memorisation; XGBoost / LightGBM / CatBoost internals; package names and code. The \emph{conceptual menu} is in scope.

